\chapter{Extracting Information from
Embedding}\label{sec-extracting-information-from-embeddings}

\textbf{Learning Objectives}

\emph{Overview starts}

\begin{itemize}
\tightlist
\item
  Understanding the concept of word Embedding
\item
  Understanding how word vectorization is used
\item
  Understanding the concept of similarity using vectors
\item
  Understanding Word2Vec and its training model
\item
  Understanding the concept of transformers
\item
  Understanding the concept of context and different embedding values
\item
  How to use word embedding to extract information
\item
  Using transformers for text summarizing
\item
  Setting up questions and Answers using transformers
\end{itemize}

\emph{Overview ends}

\section{Introduction to word
embedding}\label{introduction-to-word-embeddings}

The history of human language is tied to the evolution of our species
(Luef 2018), and it is constantly changing (Milroy 1995). Our language
encodes emotions, experiences, cultural significance, and context very
effectively. There seems to be a very tight relationship between
language and thinking, and they affect each other in several ways, such
as in the perception of time, space, or even colors (Winawer et al.
2007). The words we hear, read, or write can be transformed into vectors
of numbers which capture information not only about the word itself but
also about its context and its links to other related words. Word
embedding is a concept that has been part of both AI academic and
industrial discourse for quite some time. The name does not ring a bell
immediately, perhaps because it originates in mathematics: An
embedding is a function that maps elements from one space onto another
(Haroske and Triebel 2011). The concept has been extended to NLP, where
words, as the main components of the human language, are mapped onto a
mathematical concept represented by a continuous vector space (Mikolov, Yih, and Zweig 2013), becoming embedded there. Word embedding has
been used in research on finance, where
it is used to generate financial lexicons (Das et al. 2022),
historical corpora (Wevers and Koolen 2020), agriculture (Guo et al.
2021), and music (Kumar, Rajpal, and Rathore 2018), to name just a
few areas where it is of increasing interest.

Word embedding has seen a growing number of applications in marketing.
For example, Timoshenko and Hauser (2019) use word embedding and a
machine-learning classifier to identify customer needs from
user-generated content. Toubia, Berger, and Eliashberg (2021) use
embedding to quantify the semantic progression of a successful
(marketing) narrative. More recently, embedding has been applied to
evaluate and predict the performance of email campaigns and
recommended targeting policies for new promotions, e.g., (Ellickson et al.
2024), as well as to extract relevant contextual product characteristics and
other information from unstructured text and image data, using this
information to estimate demand or to derive product recommendations
(Sisodia, Burnap, and Kumar 2024; Compiani, Morozov, and Seiler 2023).

Word embedding can be trained on large text corpora and allows for fast
access to semantic relationships and similarities between, for example,
specific products and brand names (e.g., ``Kindle Fire'' and
``Amazon''), and also to their characteristics expressed in
words such as ``tablet'', ``e-book reader'', ``gift'', or ``customer
service''. While word embedding focuses on the process of mapping words
onto a vector space, a word vector is a directly related concept and
refers to how a word is mapped onto a numerical sequence.

\section{What is a word vector?}\label{what-is-a-word-vector}

In the following simple example, we focus on words which are related to
emotions, and we represent them through two components: polarity and
intensity. Their numerical values are encoded in a very subjective way.
Polarity, which can be negative or positive, is expressed as a range
between -1 (negative) and 1 (positive). Similarly, intensity can also
be given values over a defined range, such as (0,1). When there is no
emotional quality present in a word, the value is 0, while for intense
emotion the value is 1. The subjective emotional quality of
a word like ``joy'' can be positive and intense, and it can be
represented by two subjective numbers: {[}0.9, 0.7{]}. On the x-axis the
valence = 0.9, while on the y-axis, arousal = 0.7. See
Table~\ref{tbl-emotions} for a more complete emotional data set. In a
two-dimensional space, such as a scatterplot, and indeed in three-, four- or
n-dimensional space, there is a correspondence between points and
vectors that start at an origin and have several components, for each
dimension, one component.

\begin{longtable}[]{@{}lll@{}}
\caption{Emotional vector space}\label{tbl-emotions}\tabularnewline
\toprule\noalign{}
Word & Polarity & Intensity \\
\midrule\noalign{}
\endfirsthead
\toprule\noalign{}
Word & Polarity & Intensity \\
\midrule\noalign{}
\endhead
\bottomrule\noalign{}
\endlastfoot
Joy & 0.9 & 0.7 \\
Happiness & 0.8 & 0.6 \\
Contentment & 0.6 & 0.3 \\
Excitement & 0.7 & 0.9 \\
Anger & -0.5 & 0.8 \\
Rage & -0.9 & 1.0 \\
Sadness & -0.7 & 0.6 \\
Grief & -0.8 & 0.85 \\
Fear & -0.6 & 0.75 \\
Anxiety & -0.7 & 0.8 \\
Calmness & 0.3 & 0.2 \\
Boredom & -0.1 & 0.1 \\
\end{longtable}

While points and vectors are not strictly the same, points can represent
vectors in a vector space (see Figure~\ref{fig-instr2}). The word
``joy,'' expressed by the coordinates {[}0.9, 0.7{]}, describes
a vector and also marks its endpoint. The set of all these vectors forms
a space, in our case an emotional vector space. In real life, emotions
such as anxiety and fear are often perceived as being very similar, and this
can be reflected mathematically by their proximity on a
two-dimensional plane.

\begin{figure}

\centering{

\includegraphics{10.Chapter_files/figure-pdf/fig-instr2-1.pdf}

}

\caption{\label{fig-instr2}Words expressing emotions: Polarity
vs.~intensity \newline Authors' own image}

\end{figure}%

The distance from the origin to a point is called magnitude, but vectors
can also provide direction, so they can capture orientation as well.
While it is not possible to apply mathematics to words directly, doing so using
numerical transformations of these words is feasible, resulting in
values represented on a surface (or in a space). The similarity of these
values is equivalent to the similarity in meaning of the words,
which can be calculated through a measure called cosine similarity.

\section{Cosine similarity}\label{cosine-similarity-1}

Cosine similarity is a popular metric for evaluating the similarity between
two vectorized entities, such as vector representations of words or,
more generally speaking, text. This is accomplished by computing the
cosine value of the angle between the entities, with a value of zero (i.e.,
\(\cos(0)=1\)), indicating that the vectors are identical. If the vectors
oppose each other, they form a 180-degree angle, such as in \(2\pi\), and
\(cos(2\pi)=-1\), indicating that they are opposite. If the vectors are
perpendicular, the angle is 90 degrees, which can be expressed as
\(\frac{\pi}{2}\), and \(cos(\frac{\pi}{2})=0\), indicating
they are unrelated. In case you are rusty on this topic, a refresher on
trigonometry can be useful here, especially of the trigonometric circle.

To better illustrate this with an example, let us consider two words,
``fear'' and ``anxiety,'' along with their vector representations: fear
{[}0.25, 0.7{]} and anxiety {[}0.25, 0.8{]}. Let us denominate the angle
between these two vectors by \(\theta\). The cosine similarity between
the vectors \(\mathbf{F}\) and \(\mathbf{A}\) is given by the formula:

\[
\cos(\theta) = \frac{\mathbf{F} \cdot \mathbf{A}}{|\mathbf{F}| \times |\mathbf{A}|}
\]

The numerator represents the dot product of the vectors:
\(\textbf{F} \cdot \textbf{A} = (-0.6 \times -0.7) + (0.7 \times 0.8)\),
while the denominator is the product of the magnitudes of the vectors:
\textbar{}\textbf{F}\textbar{} \times \textbar{}\textbf{A}\textbar{} =
\(\sqrt{(-0.6)^2 + (0.7)^2} \times \sqrt{(-0.7)^2 + (0.8)^2}\).

Plugging these values into the formula, we obtain a cosine similarity
value of approximately 0.999948. This high degree of similarity (close to 1)
shows that these emotional expressions are perceived as being very similar,
both in everyday understanding and in their numerical representation. A
cosine similarity close to 1 indicates that there is a very small angle
between these two vectors.

Instead of using two features, we can generalize this into many, see
Table~\ref{tbl-emot2}. We started with a basic example in two dimensions
using polarity and intensity, and we have extended the example by adding
three more dimensions to better capture the spectrum of the emotional characteristics of the
words. In addition, Arousal deals with the state of alertness or the energy of an emotion, the valence represents how pleasant or
unpleasant an emotion is, and dominance shows the how much control or
influence an emotion exerts. While visualizing in more than two (or
three) dimensions is tricky, this is not very
demanding computationally. If we take a look at only the numbers, they may seem
meaningless at first glance, but they are crucial in computational
linguistics, as we can carry out several operations using these vectors and
land in a region, i.e., reach a point in the vector space where the
combination of features and words will show major densities and the positions of
the vectors in nearby locations. Let us take a closer look at the row
elements that are similar in the vector space, are they similar in
meaning too?

\begin{longtable}[]{@{}llllll@{}}
\caption{Emotion characteristics table}\label{tbl-emot2}\tabularnewline
\toprule\noalign{}
Word & Polarity & Intensity & Arousal & Valence & Dominance \\
\midrule\noalign{}
\endfirsthead
\toprule\noalign{}
Word & Polarity & Intensity & Arousal & Valence & Dominance \\
\midrule\noalign{}
\endhead
\bottomrule\noalign{}
\endlastfoot
Joy & 0.9 & 0.7 & 0.8 & 0.9 & 0.7 \\
Happiness & 0.8 & 0.6 & 0.7 & 0.85 & 0.7 \\
Contentment & 0.6 & 0.3 & 0.4 & 0.75 & 0.5 \\
Excitement & 0.7 & 0.9 & 0.9 & 0.85 & 0.8 \\
Anger & -0.5 & 0.75 & 0.8 & 0.15 & 0.4 \\
Rage & -0.9 & 0.95 & 0.95 & 0.1 & 0.3 \\
Sadness & -0.7 & 0.6 & 0.5 & 0.35 & 0.4 \\
Grief & -0.8 & 0.85 & 0.7 & 0.2 & 0.3 \\
\textbf{\emph{Fear}} & \textbf{-0.6} & \textbf{0.7} & \textbf{0.9} &
\textbf{0.3} & \textbf{0.4} \\
\textbf{\emph{Anxiety}} & \textbf{-0.7} & \textbf{0.8} & \textbf{0.9} &
\textbf{0.3} & \textbf{0.4} \\
Calmness & 0.3 & 0.2 & 0.3 & 0.65 & 0.5 \\
Boredom & -0.1 & 0.1 & 0.2 & 0.5 & 0.35 \\
\end{longtable}

Notice that the two emotions, anxiety and fear, are expressed as
vectors. If you take the numbers in the corresponding positions, they
are very similar. We can also notice a fundamental insight: Similarity in
numbers equals similarity in significance. Because we have vectors, we can
calculate analogies, following the classic example, king - man + woman =
queen (Mikolov et al. 2013). We can, therefore, check: if anxiety stands next to
fear, joy stands next to what? If you do the computations, you obtain a
vector that is very close to happiness. This is the way word
embedding can be used to identify analogies. In
Figure~\ref{fig-rect1} we can see how the word embedding takes the form of a matrix (the pictures that are not
otherwise stated as authors' are inspired by Alammar and Grootendorst
2024). The features, such as polarity, intensity, arousal, valence,
and dominance from Table~\ref{tbl-emot2}-, and others, appear vertically
in the matrix, while words from the vocabulary, such as joy,
happiness, and many more, are seen in the rows.

\begin{figure}

\centering{

\includegraphics[width=0.3\textwidth,height=\textheight]{figure/chapter10/rect1.png}

}

\caption{\label{fig-rect1}Embeddings matrix}

\end{figure}%

Now let us try to understand how these numbers are generated. An
intuitive lead is given by the fact that, in a large corpus, similar
words appear in similar contexts through their co-occurrence and
generate similar vectors, i.e., similar embedding. Word embedding is
used to assist or to improve tasks you have already seen, such as sentiment
analysis, topic modeling, classification, and even translation, by
using the semantic information of words encoded in the vectors. A
popular algorithm to establish word embedding that tries to capture
semantic similarities between words is \texttt{Word2Vec} (Mikolov
et al. 2013).

\section{Creating word embedding}\label{word-embeddings-creation}

The work of linguists such as Ferdinand de Saussure (1857–1913), or
anthropologists who made strong contributions to linguistics, such as Franz
Boas (1858–1942), underlined the importance of context in
understanding language. In 1957, in his distributional theory which studies how the distribution of
words in a text can reveal their meaning, the British linguist John Rupert Firth coined the phrase ``[y]ou shall
know a word by the company it keeps'' (see Evert et al. 2008). The simple
intuition behind this concept is that the meaning of the words is
determined by the context they appear in, i.e., by the other words
surrounding them. Some semantic aspects of the words can be defined by
homonymy or polysemy. Homonymy means that words can look
the same but have unrelated meanings. Terms like ``mean'' can ``mean''
(pun intended) or signify unpleasant, primitive, or average. Polysemy
describes situations where words can have related meanings, as in the
sentence ``the frame was beautiful, but the painter was framed'':The
first occurrence of the word ``frame'' refers to the frame of a
painting, while the second one refers to the painter being falsely
incriminated. In both homonymy and polysemy, context analysis is
necessary to disambiguate the target words.

An NLP algorithm, following Firth's theories and focusing on
words in context but which does not capture polysemy or homonymy directly is
\textbf{Word2Vec}. As with the example of the contextualized
emotional vector space, it represents words as a sequence of
numerical values in a more generic vector space. Note that the same word
in \textbf{Word2Vec}, despite having a different context, cannot be represented
by different sequences of numbers. There are two main procedures
performed by \textbf{Word2Vec}: continuous bag-of-words (CBOW) and skip-gram.
The first one predicts what a target word will be based on the surrounding words, and
the second tries to predict the surrounding words based on a target
word. We will explain how this is accomplished. Before
doing so, however, we need you to familiarize yourself with some more basic
concepts.

The words in a corpus will be transformed into embedded numbers created
through the training of a neural network. Like all the tasks in
artificial intelligence, the purpose of this operation is predictive. Let us consider Figure~\ref{fig-w2vconcept} and understand how
\textbf{Word2Vec} is implemented conceptually.

\begin{figure}

\centering{

\includegraphics[width=0.6\textwidth,height=\textheight]{figure/chapter10/w2vconcept.png}

}

\caption{\label{fig-w2vconcept}The Word2Vec operation scheme}

\end{figure}%

Initially, we have the vocabulary, and each word in the vocabulary is
represented as a one-hot vector. This representation, similar to
Figure~\ref{fig-hotencoding}, considers all the words as vectors of the
size of the vocabulary. When we select a specific word, that word will
have a one corresponding to the index of the word, and all the other
ones are zero. It is like focusing on the existence of that word, see
Table~\ref{tbl-onehot}. This way, we map each word to a position vector.
The words do not necessarily have to be in alphabetical order, they can also be in
order of appearance. If a word is repeated, it takes the index of the
first occurrence.

\begin{longtable}[]{@{}lll@{}}
\caption{One-Hot Encoding of a
sentence}\label{tbl-onehot}\tabularnewline
\toprule\noalign{}
Word & Index & One-Hot Vector \\
\midrule\noalign{}
\endfirsthead
\toprule\noalign{}
Word & Index & One-Hot Vector \\
\midrule\noalign{}
\endhead
\bottomrule\noalign{}
\endlastfoot
The & 1 & {[}1, 0, 0, 0, 0, 0, 0, 0, 0, 0, 0, 0, 0{]} \\
summer & 2 & {[}0, 1, 0, 0, 0, 0, 0, 0, 0, 0, 0, 0, 0{]} \\
sale & 3 & {[}0, 0, 1, 0, 0, 0, 0, 0, 0, 0, 0, 0, 0{]} \\
campaign & 4 & {[}0, 0, 0, 1, 0, 0, 0, 0, 0, 0, 0, 0, 0{]} \\
went & 5 & {[}0, 0, 0, 0, 1, 0, 0, 0, 0, 0, 0, 0, 0{]} \\
viral & 6 & {[}0, 0, 0, 0, 0, 1, 0, 0, 0, 0, 0, 0, 0{]} \\
on & 7 & {[}0, 0, 0, 0, 0, 0, 1, 0, 0, 0, 0, 0, 0{]} \\
social & 8 & {[}0, 0, 0, 0, 0, 0, 0, 1, 0, 0, 0, 0, 0{]} \\
media & 9 & {[}0, 0, 0, 0, 0, 0, 0, 0, 1, 0, 0, 0, 0{]} \\
and & 10 & {[}0, 0, 0, 0, 0, 0, 0, 0, 0, 1, 0, 0, 0{]} \\
increased & 11 & {[}0, 0, 0, 0, 0, 0, 0, 0, 0, 0, 1, 0, 0{]} \\
brand & 12 & {[}0, 0, 0, 0, 0, 0, 0, 0, 0, 0, 0, 1, 0{]} \\
awareness & 13 & {[}0, 0, 0, 0, 0, 0, 0, 0, 0, 0, 0, 0, 1{]} \\
\end{longtable}

Now having represented the vocabulary, we can take a single input word, for
example ``summer'', which has been fed into the neural network (this is the skip-gram
approach, see Figure~\ref{fig-skipgram}). In this image the vocabulary
has 10,000 unique words. The hidden layer is considered a lookup table
which is nothing but the word-vector representation of the word
``summer''. This vector is then sent to the output layer, which is a
classifier. The projection layer, where the embedding matrix resides,
is used as the learned mapping which transforms the one-hot vector into a
continuous vector (i.e., embedding). The softmax function converts the output
into a probability distribution over the vocabulary. The result is a
vector of the same dimension as the vocabulary but containing the
probabilities of which words are close to the input word. For more
information, see McCormick (2016).

\begin{figure}

\centering{

\includegraphics[width=0.85\textwidth,height=\textheight]{figure/chapter10/skip_gram_net_arch2.png}

}

\caption{\label{fig-skipgram}Skipgram architecture. (Picture inspired by
and edited from McCormick.com)}

\end{figure}%

Here we presented only the kkipgram model, but alternatively the CBOW
approach can be used too, where you can take multiple context words
and use them as input, often by averaging their one-hot vectors
(see, Q. Wang et al. 2017).

As a real-world example, when you are writing your favorite instant-messaging app uses an
algorithm that tries to guess what your next word will be,
based on what you have already typed. Let us assume that you want
to send the sentence we presented in Table~\ref{tbl-onehot} to your
business partner: ``The summer sale campaign went viral on social media
and increased brand awareness''. Which words would a trained model, such
as one created with \textbf{Word2Vec}, predict for you to type after
``increased''? Would it be sales, traffic, engagement, or something
else? Figure~\ref{fig-trained} represents a stylized version of an
embedding, where the trained language model represents the weights of
the neural network, which would propose a series of ``next words'' using a
specific probability.

\begin{figure}

\centering{

\includegraphics[width=0.8\textwidth,height=\textheight]{figure/chapter10/trained.png}

}

\caption{\label{fig-trained}A trained model for guessing the next word}

\end{figure}%

The two words in the input, ``and'' and ``increased'', are the tokens we
want the model to process (see the left-hand side of
Figure~\ref{fig-trained}) and are mapped onto their corresponding
embedding, which is the output of a trained Word2Vec model. This
embedding is shown in the embedding matrix, where each row
corresponds to a word in the vocabulary, such as ``a'', ``aa'',
``aardvark'', ``aardwolf'', etc., and each column corresponds to one
feature of the embedding vector. The neural network computes a score for
each word in the vocabulary, which is then converted into a probability
distribution using a function such as softmax. For example, given the
input words provided, the model assigns probabilities to candidates for the next
word, such as ``sales'' ( with a probability of 0.36), ``brand'' (0.38), and
others, based on how well they fit the patterns it has learned. If the model
has learned that through training  that a phrase like ``and increased sales''
is common, the probability of seeing ``sales'' after ``and increased'' is
high. Next, we will see how these models can be trained.

\section{Algorithms in language model
training}\label{algorithms-in-language-model-training}

For training word embedding, we first need to provide the training
algorithm with the context of the surrounding individual words. This is
accomplished by defining the concept of a sliding window. Moving the
window over the text generates several context-target pairs. They
capture the local contexts of words by building structured examples,
helping the model to learn which words appear together frequently and
identifying meaningful relationships inside the corpus. The window
represents a certain number of words before (and after) a target word
and can generate a data set to train a model. There can be any number of words, but some capture information better than others. This
information assesses which words are near other words. Statistically
speaking, if a word appears in a specific context, other words related
to that context will be there too. Let us start with a simple example,
with a window of two words, for the sentence we have already used: ``The
summer sale campaign went viral on social media and increased brand
awareness''.

\begin{figure}

\centering{

\includegraphics[width=0.8\textwidth,height=\textheight]{figure/chapter10/sliding_window_table.png}

}

\caption{\label{fig-training}The sliding window concept}

\end{figure}%

In Figure~\ref{fig-training}, two tables can be seen. Each row shows how the sliding window, which is made of two words, advances. The word
after that is the target word. Visually, the two-word window slides to
the right, capturing just a part of the sentence. A data set is built in the table to the right to save all the information captured from the
sliding window. Each pair of words used as an input will have a
single word associated with it as an output of the data set we have generated.

To take this a step further and to morph the window differently, we can select two
words before and two words after a target word, and use the basic
idea of CBOW (Jang, Kim, and Kim 2019). So, given an even
number of words as an input defined by a specific window size (for example
two), the output is the word in the middle. Note that in this approach
the target word is not part of the input, but it is the word that is
omitted and then predicted based on the surrounding context.

\begin{figure}

\centering{

\includegraphics[width=0.8\textwidth,height=\textheight]{figure/chapter10/cbow.png}

}

\caption{\label{fig-cbow}CBOW approach}

\end{figure}%

In Figure~\ref{fig-cbow}, the sliding window technique is applied to a sentence using CBOW. Each line in the table is a
subset of the sentence, showing how the window captures different
segments as it moves from left to right. Here we are trying to predict a target
word based on the surrounding context words, two ahead and two before it.
The two shades of gray represent the sliding window words (the input in
light gray) and the target word (the output in dark gray).

A complementary approach is to have an input word and guess what the two
words before and the two words after it would be. This is demonstrated
in Figure~\ref{fig-trainingprocess}, which adds further information and
shows how the training process for the word-embedding model process is
carried out. For example, by applying the sliding-window approach above
(i.e., using one input and two candidate target words before and after
the input), the input word ``sale'' is associated with ``the'',
``summer'', ``campaign'', and ``went'', and the training dataset
continues with the next input word and its related words. At the very
beginning of the process, the model contains no information provided by
random embedding, meaning that the model is untrained at the outset. If you have the word
``sale'', which word would be a neighbor? In an untrained model, the
prediction would be random, and, likely, it could not learn any meaningful
word representations (yet); its predictions would be totally inaccurate.

\begin{figure}

\centering{

\includegraphics[width=0.8\textwidth,height=\textheight]{figure/chapter10/training_process.png}

}

\caption{\label{fig-trainingprocess}Word-embedding training process}

\end{figure}%

In Figure~\ref{fig-trainingprocess}, the ''Model Prediction'' column shows
the probabilities the model assigns to each word in the vocabulary as
a possible neighboring word for ``sale''. The ``Actual Target'' column
shows an example of the correct neighboring word as a so-called
``one-hot encoded'' vector (in this case, ``the''), which means applying
a value of 1 to the correct word and a zero to all the others. The ''Error''
column represents the difference between the actual target and the value from the untrained model. This error value is used to update
the model's parameters through backpropagation and retraining (Goh
1995; Rumelhart et al. 2013). Backpropagation is a fundamental method
used to train neural networks. The error value is ``sent'' backwards
through the network to adjust the weights. The operation is performed
several times until the network learns to give more accurate
predictions.

After we had done the same process to cover all the samples in the data set,
we finished one epoch of training. The process is repeated for several
epochs until the model becomes stable and yields significant results, or
until a certain planned number of epochs have been performed. Eventually,
we will have acquired a trained model, and we can extract the word-embedding
matrix from it, which is our main goal.

\section{Optimizing embedding}\label{optimizing-embeddings}

In calculating how to predict the next word, the most computationally
expensive part is the projection of the output vocabulary. To improve
this task, first, we need to generate high-quality embedding; then, we
need to use this embedding to train a language model (to carry out
next-word prediction). At this point, we assume that we have a fully trained embedding, and we can now switch from the
training phase to prediction. Now, instead of
predicting a neighboring word directly, the model calculates a score for
a given input-output word pair to indicate if they are neighbors or not
(1 or 0). This switch effectively changes the model from a regular
neural network for next-word prediction into a logistic-regression
model, which is simpler and faster to compute. The structure of our
data set changes slightly through the addition of a new column assessing the presence
of two neighboring words. We do not now have to look to guess the next
word anymore, or if the words are neighbors or not. This operation can
be done very quickly, but with one caveat.:If all the examples are positive,
meaning that each target value is 1, we have a model that will always return
1, meaning it has perfect accuracy, which is a clear case of overfitting.
To overcome this, negative samples of words which are not neighbors are
sampled at random from the vocabulary, and a target value of 0 is assigned
to those word pairs. This concept was introduced by Gutmann and Hyvärinen
(2010) and is shown in Figure~\ref{fig-neg}.

\begin{figure}

\centering{

\includegraphics[width=0.8\textwidth,height=\textheight]{figure/chapter10/neg.png}

}

\caption{\label{fig-neg}Skipgram with negative sampling}

\end{figure}%

Before the training process starts, we have to check which words are
involved and how many words are in the vocabulary of the corpus. Let
us combine the skipgram and negative-sampling methodologies and show how
\textbf{Word2Vec} uses them during the training process. At the start of the
training phase, it creates the the embedding and context matrices, see
Figure~\ref{fig-embeddingcontext}. You can see that they are
exactly the same size, containing the same vocabulary and the same
number of features. What differs between them, is that the values of
the features will be different, and after the conclusion of the training
process, there will be two distinct matrices, but of the same size.

\begin{figure}

\centering{

\includegraphics[width=0.8\textwidth,height=\textheight]{figure/chapter10/embedding_context.png}

}

\caption{\label{fig-embeddingcontext}Embedding and context matrices}

\end{figure}%

As in any training process, the matrices are initially populated with
random values. Both have unique embedding sequences for each word in the
vocabulary, meaning that a word in the embedding matrix will have
different values in the context matrix. Each word is represented as a
row. The second dimension of the matrices holds the number of features,
a concept similar to the emotional space where we had just two: valence
and intensity. While 300 is a commonly employed value, the number of
features can range from between 50 to 500, or any number you find suitable.
It is important to note that they do not have any predetermined label
assigned to them. Figure~\ref{fig-finaltwo} shows how the algorithm
updates the model's parameters. A series of computational steps are performed
to improve the embedding of the words, i.e., the correspondence between the words
and their context.

\begin{figure}

\centering{

\includegraphics[width=0.8\textwidth,height=\textheight]{figure/chapter10/finaltwo.png}

}

\caption{\label{fig-finaltwo}The Word2vec training process}

\end{figure}%

The training mechanism of the two matrices, embedding and context,
follows the standard procedure of training a neural network using
backpropagation mentioned in the previous section. This entails a forward
pass, a loss or error function calculation, a backward pass (called
backpropagation), and a parameter update. At the beginning, during the first
forward pass, the values in the matrices, also called weights or
parameters, are populated with randomly assigned values to facilitate
effective learning. If we look at Figure~\ref{fig-finaltwo} again, the
dot product of input and output, \(\text{input}\cdot\text{output}\), is
calculated using random values. The result is passed on to the sigmoid()
function and yields a value between 0 and 1, which, in turn, is compared
to the value in the \textbf{target} column to compute the error, hence, carrying out the loss function. This is an extremely
important value because the model parameters are
iteratively updated based on this to make the error value as close as possible to 0.
The updating procedure is done through backpropagation {[}reference{]}
from the output layer to the input, updating the parameters through an
optimization algorithm called gradient descent.

\subsubsection{Details of the training process}\label{training-process-details}

Training a neural network, as we are doing here, can involve
optimizing the process using hyperparameters. This term is a combination of
two words: ``hyper'', from Greek, meaning ``above'' or ``beyond'', and
parameter, which was initially composed of two words, ``para'', meaning
``beside'', and ``metron'', meaning ``measure''. So, a hyperparameter is
a variable which operates before training begins at a higher level and
can influence the final model. One such parameter is the window size,
which can have a strong impact on the outcome. A small window learns
specific and tight relationships between words, such as ``milk'' and
``shake''. A larger window learns broader relationships, such as ``milk''
and ``animal''. Another parameter is the dimensionality of the word
vector (the columns in Figure~\ref{fig-rect1}). A higher dimension
captures subtle relationships between words, but it can lead to
overfitting and increased computational costs. A lower dimension is
faster but might not capture all the nuances of the relationships
between words.

At this point we are calculating the gradient of the error of the embedding vectors. This involves
computing the partial derivatives of the loss function for
each element of the embedding and context vectors. This step is
essential to update the parameters, both for embedding and the context
matrices. The generic updating formula can be expressed as:

\[
\theta = \theta - \eta \cdot \nabla_{\theta} J(\theta)
\] where \(\theta\) is the parameter being updated, \(\eta\) is the
learning rate, \(\nabla_{\theta} J(\theta)\) is the gradient of the cost
function \(J(\theta)\) for the parameters \(\theta\).

Both the embedding and context matrices follow the same steps. Each
word's vector is updated by slowly moving in the direction that reduces
the error. This involves subtracting a constant fraction, identified by
\(\eta\), of the gradient of the loss from the current vector for the embedding
vector.

Let us denote the embedding vector of the input word as \(v\) and the
context vector of the output word as \(u\), and the error calculated as
E. The updates can be expressed mathematically as:

\[
v = v - \eta \cdot \frac{\partial E}{\partial v}
\]

\[
u = u - \eta \cdot \frac{\partial E}{\partial u}
\]

These equations represent gradient descent updates for the parameters
\(v\) and \(u\), where \(\eta\) is the learning rate, and
\(\frac{\partial E}{\partial v}\) and \(\frac{\partial E}{\partial u}\)
are the gradients of the error function E for \(v\) and
\(u\).

The partial derivatives of the error function E for the
variables \(v\) and \(u\), \(\frac{\partial E}{\partial v}\) and
\(\frac{\partial E}{\partial u}\), depend on the specific form of the
error function used. It often involves the sigmoid function of the dot
product of \(v\) and \(u\), compared against the actual target (1 or 0
in our case). This process is repeated for each training example,
gradually modifying the values of the embedding by minimizing the
overall error in predicting the context words given the input words. The
loss function J and the error (E) are closely related, and in many cases,
they are essentially the same thing. The loss function (J) is typically
defined for an entire dataset. It represents the overall performance of
the model across multiple examples, while the E usually refers
to the error of a single training example. In the context of Word2Vec
and similar models, we often use the terms interchangeably for a single
example. The loss function J for the entire dataset would be the sum (or
average) of the individual errors E for each training example.

Let us focus on a practical and simple example to understand exactly
what happens. Following how the procedure was described in
Figure~\ref{fig-finaltwo} we can define our error function as:

\[
E = t - \sigma(v \cdot u)
\] where \(\sigma()\) is the sigmoid function, and \(t\) is the target
(0 or 1).

Calculating the partial derivatives, knowing that the derivative of the
sigmoid function is \(\sigma'(x) = \sigma(x)\left(1 - \sigma(x)\right)\)
and the derivative of a constant value, such as t, results in \(t'=0\).

\[
\frac{\partial E}{\partial v} = -\sigma(v \cdot u)(1 - \sigma(v \cdot u)) \cdot u
\]

\[
\frac{\partial E}{\partial u} = -\sigma(v \cdot u)(1 - \sigma(v \cdot u)) \cdot v
\]

We can then update the vectors to: \[
v_{\text{new}} = v_{\text{old}} - \eta \cdot \left(-\sigma(v_{\text{old}} \cdot u_{\text{old}})(1 - \sigma(v_{\text{old}} \cdot u_{\text{old}})) \cdot u_{\text{old}}\right)
\]

\[
u_{\text{new}} = u_{\text{old}} - \eta \cdot \left(-\sigma(v_{\text{old}} \cdot u_{\text{old}})(1 - \sigma(v_{\text{old}} \cdot u_{\text{old}})) \cdot v_{\text{old}}\right)
\]

We can use a numerical example to better understand the whole process.

\[
v = \begin{bmatrix} 0.5, 0.3 \end{bmatrix}
\] \[
u = \begin{bmatrix} 0.1, 0.2 \end{bmatrix}
\] \[
t = 1 \quad \text{(positive example)}
\] \[
\eta = 0.1 \quad \text{(learning rate)}
\]

Calculate: \[
v \cdot u = 0.5 \cdot 0.1 + 0.3 \cdot 0.2 = 0.11
\] \[
\sigma(0.11) \approx 0.5275
\] Error: \[
E = 1 - 0.5275 = 0.4725
\]

Gradient for (v): \[
-\sigma(0.11) \cdot (1 - \sigma(0.11)) \cdot \begin{bmatrix} 0.1, 0.2 \end{bmatrix} \approx \begin{bmatrix} -0.0249, -0.0498 \end{bmatrix}
\]

Update (v): \[
v_{\text{new}} \approx \begin{bmatrix} 0.5, 0.3 \end{bmatrix} - 0.1 \cdot \begin{bmatrix} -0.0249, -0.0498 \end{bmatrix} \approx \begin{bmatrix} 0.50249, 0.30498 \end{bmatrix}
\] You can see that the new values of the vector are close to the
original ones, but by performing this task several times, the values
will be adjusted to obtain an optimal error function.

\section{Other types of embedding and models}\label{other-embeddings-and-models}

There are other popular algorithms for representing the vector space, such
as Global Vectors (GloVe; Pennington, Socher, and Manning 2014),
developed at Stanford, and FastText (Bojanowski et al. 2017), developed
by Facebook (now Meta). The aim is, as with Word2Vec, to capture
meaning in the text using an unsupervised algorithm. GloVe is based
on the concept of the co-occurrence matrix, which is built for an entire
corpus, summarizing how often word pairs occur together. The main
idea is that ratios of co-occurrence probabilities contain useful
information. Words like ``coffee'' and ``cappuccino'' would have similar
vector representations, as they frequently appear together (co-occur) in
certain contexts: social meetings, breakfasts, in bars, etc. While
Word2Vec focuses on context using sliding windows, GloVe uses a
co-occurrence matrix for the entire corpus. FastText, instead, builds on Word2Vec, incorporating sub-word
information to improve their vector representation, meaning that words
are broken down into character n-grams. This improves the
probability of including underrepresented words which are
out-of-vocabulary. Such examples can be words which do not appear in the
training session but appear in a corpus and should be analyzed. So, if a new
word appears, it is broken down into familiar sub-words (or n-grams).
FastText performs well at capturing how morphological details are defined.
Languages with many words can benefit the most from using this algorithm. For
example, German is a good fit because it has many compound words and
inflectional changes, and it is morphologically rich.

There are an increasing number of other methods which are designed to
learn word embedding. For example, Doc2Vec (Le and Mikolov 2014)
associates a unique vector with each document or paragraph, focusing on
larger portions of text rather than single words. Sent2Vec (Pagliardini,
Gupta, and Jaggi 2017) builds sentence-level representations from
uni-grams and bi-grams, Embedding from language models (ELMo; Peters et
al. 2018) focuses on the context of the sentence to generate word vectors, while bidirectional encoder representations from transformers (BERT; Devlin
2018) uses the transformer architecture to generate context-aware
embedding, or GPT (Radford 2018), another transformer-based
architecture, which became extremely popular, starting with version GPT-3 in 2020 (Brown 2020). Transformers represent a further
development of word vectors as they allow for the capturing of context, or
generating text that fits nicely into a specific context. Below we
present a brief discussion of the building blocks of this architecture.

\subsection{Transformers}\label{transformers}

A transformer is an important type of neural network. It uses word
embedding to learn context and implicit meaning by detecting
relationships between words in sentences. Transformers are the main
architectural framework for the unsupervised learning models underlying BERT or
GPT. The link to word embedding discussed in the previous section is that the
input tokens, which are words or sub-words, are transformed into arrays
of numbers, but with a twist. Unlike in Word2Vec, transformers add a
position encoding to the vectors to ensure that the model captures the
position of each word in a sentence. This dual information is passed on
to the transformer model. They rely on a self-attention
mechanism (introduced by Vaswani 2017) to process input sequences in
parallel. The transformer architecture includes an encoder,which
processes the input sequence, and sometimes, but not always, a decoder
which generates the output sequence. In addition, there are multi-head
self-attention layers, which run in parallel and focus on different
elements of the input. After establishing an attention layer, it applies feedforward layers, and then uses positional encoding. There are three types of transformer: (i) encoder-only, which are designed to understand
text (e.g., BERT); (ii) decoder-only, which are designed to generate
text (e.g., GPT); and (iii) encoder-decoder sequences, which are useful
for translating languages. This section drew on the basic explanation of
transformers by Alammar and Grootendorst (2024).

\subsubsection{Encoder-decoder}\label{encoder-decoder}

To illustrate and to better understand the transformer concept,
Figure~\ref{fig-transformer} presents a encoder-decoder
architecture. The diagram shows how language translation can be
processed. The phrase ``to be happy'' is sent to an encoder made up of six
layers, each representing a state of processing. The output from the
decoder leads to the correct translation, in this case the German ``glücklich
sein''.

\begin{figure}

\centering{

\includegraphics[width=0.8\textwidth,height=\textheight]{figure/chapter10/transformer.png}

}

\caption{\label{fig-transformer}Transformer, encoder and decoder}

\end{figure}%

A more detailed representation of the encoder and decoder layers sheds
more light on their internal working (see
Figure~\ref{fig-encoderanddecoder}). The encoder consists of a feed-forward neural network which processes data through a series of
transformations, and a self-attention mechanism to evaluate the
importance of different sections of the input, which aims to capture
dependencies and relationships. The decoder has its own
feed-forward neural network, used to produce the output, an
encoder-decoder attention element, which allows the decoder to focus on the
important parts of the input sequence, and a self-attention mechanism
for considering the previous output when generating the next token.

\begin{figure}

\centering{

\includegraphics[width=0.8\textwidth,height=\textheight]{figure/chapter10/encoder_decoder.png}

}

\caption{\label{fig-encoderanddecoder}Transformer: encoder and decoder
detail}

\end{figure}%

To illustrate, let us now focus on the encoder and add more detail
by zooming in on one layer. Let us take the vectorized
sequence (as in Word2Vec) for the words ``To'', ``be'', and
``happy'', each denoted in turn as \(x_1, x_2,x_3\). When applying the
self-attention mechanism, we receive another embedding, which
captures the context. For example, ``to'' is influenced by ``be'' and
``happy'',  also it might consider the relative importance of the words, their position, or some distinct features
of the words may be emphasized, or the semantic relationships
between words highlighted. The new embedding, shown by \(z_1, z_2, z_3\), is a more
refined and context-aware representation. The transformation
of \(z_1, z_2, z_3\) into \(r_1, r_2, r_3\) within the encoder layers
is achieved by applying a point-by-point feed-forward network to
process the attention output further and extract and refine features from the
z vectors (for more detailed information, see Vaswani 2017). Including many layers is important so that each of them can refine the
vector representation for greater precision.

\begin{figure}[H]

{\centering \includegraphics[width=0.8\textwidth,height=\textheight]{figure/chapter10/encoder.png}

}

\caption{A detailed view of an encoder}

\end{figure}%

After processing all the layers, the final output of the encoder
stack will contain a context-aware representation of the input sequence
and serves as the bridge between the encoder and decoder. The main idea
is that while a word can be translated accurately without such
context-rich information, doing the same for a phrase would result in much greater precision if the context were included. The resultof the
encoding of the phrase ``to be happy'' is a set of contextualized
representations for each word in the sentence. While as an
input we have ``to be happy'' (3 tokens), the output contains three
vectors, one for each word, typically of the dimension \(d_{model}=512\) or
\(d_{model}=768\), as in many examples.

The vector for ``to'' becomes:
\(\vec{v}_{to} = [v_{to,1}, v_{to,2}, \ldots, v_{to,d_{model}}]\), the
vector for ``be'' becomes:
\(\vec{v}_{be} = [ v_{be,1}, v_{be,2}, \ldots, v_{be,d_{model}}]\), and
the vector for ``happy'' becomes:
\(\vec{v}_{happy} = [ v_{happy,1}, v_{happy,2}, \ldots, v_{happy,d_{model}}]\).

In other words, each of the vectors
\(\vec{v}_{to}, \vec{v}_{be}, \vec{v}_{happy}\) contains information not
only about the word itself, but also about its context inside the
sentence. Using this layered approach, the self-attention mechanism provides
a result which captures the interactions between words, and the feed-forward
layers in the encoder refine these contextual representations. The
numerical representation expressed as a vector for the word ``to''
contains information about its relationship to ``be'' and ``happy'', the
one for ``be'' contains information about its relationship to ``to'' and
``happy'', and the one for ``happy'' contains information about its
relationship to ``to'' and ``be''. These final output vectors capture
semantic and syntactic information about each word in the context of the
whole sentence. Transformers can be used to carry out several tasks,
such as classification, translation, or even language generation. But
they have also potential for being employed for other (predictive or
generative) tasks involving sequential information other than language,
such as for creating histories of customer-firm interactions (e.g., Kim, McCarthy,
and Lee 2025) or market basket data (Gabel and Ringel 2024).

In the next section, we will focus on the more basic Word2Vec
application, before going on to use transformers.

\section{Word2vec in practice}\label{word2vec-in-practice}

Perhaps the best way to understand how word embedding works is to use a
practical example. One of the most common tasks that Word2Vec is used for is
to find analogies which prove the relationship between words. To
illustrate this, we have used our book's data set and have selected the Content and Brand columns. The purpose is to check which words
of interest are embedded in similar vectors which capture similar
contexts.

\begin{Shaded}
\begin{Highlighting}[]
\FunctionTok{library}\NormalTok{(data.table)}
\FunctionTok{library}\NormalTok{(lubridate)}
\FunctionTok{library}\NormalTok{(word2vec)}

\CommentTok{\# Read the dataset into a data.table}
\NormalTok{prodreviews }\OtherTok{\textless{}{-}} \FunctionTok{fread}\NormalTok{(}\StringTok{"data/Product Reviews.csv"}\NormalTok{, }
                     \AttributeTok{select =} \FunctionTok{c}\NormalTok{(}\StringTok{"Content"}\NormalTok{, }\StringTok{"Brand"}\NormalTok{))}
\end{Highlighting}
\end{Shaded}

The \texttt{word2vec()} function has transformed and embedded the words using the CBOW method. We have selected a dimension of 200, a
sliding window of 5 words, and limited the number of iterations to 30.

\begin{Shaded}
\begin{Highlighting}[]
\CommentTok{\# Train word2vec model}
\NormalTok{model }\OtherTok{\textless{}{-}} \FunctionTok{word2vec}\NormalTok{(}\AttributeTok{x =} \FunctionTok{tolower}\NormalTok{(prodreviews}\SpecialCharTok{$}\NormalTok{Content), }\AttributeTok{type =} \StringTok{"cbow"}\NormalTok{, }\AttributeTok{dim =} \DecValTok{200}\NormalTok{, }\AttributeTok{window=}\DecValTok{5}\NormalTok{, }\AttributeTok{iter =} \DecValTok{30}\NormalTok{)}
\end{Highlighting}
\end{Shaded}

As a basic operation we can search for terms such as ``apple'' and
``samsung'', and find which words are similar (i.e., close to the focal
brand name) in the word-vector space.

\begin{Shaded}
\begin{Highlighting}[]
\CommentTok{\# Find similar words to apple}
\NormalTok{result\_a }\OtherTok{\textless{}{-}} \FunctionTok{predict}\NormalTok{(model, }\FunctionTok{c}\NormalTok{(}\StringTok{"apple"}\NormalTok{), }\AttributeTok{type =} \StringTok{"nearest"}\NormalTok{, }\AttributeTok{top\_n =} \DecValTok{5}\NormalTok{)}
\FunctionTok{print}\NormalTok{(result\_a)}
\end{Highlighting}
\end{Shaded}

\begin{verbatim}
$apple
  term1   term2 similarity rank
1 apple  apples  0.7145024    1
2 apple     ios  0.7050035    2
3 apple samsung  0.6755056    3
4 apple     rim  0.6742749    4
5 apple      hp  0.6641442    5
\end{verbatim}

The closest words to ''apple'' include comparisons to Samsung, Research in Motion (rim), HP and a reference to its operating
system, iOS.

\begin{Shaded}
\begin{Highlighting}[]
\CommentTok{\# Find similar words to samsung}
\NormalTok{result\_s }\OtherTok{\textless{}{-}} \FunctionTok{predict}\NormalTok{(model, }\FunctionTok{c}\NormalTok{(}\StringTok{"samsung"}\NormalTok{), }\AttributeTok{type =} \StringTok{"nearest"}\NormalTok{, }\AttributeTok{top\_n =} \DecValTok{5}\NormalTok{)}
\FunctionTok{print}\NormalTok{(result\_s)}
\end{Highlighting}
\end{Shaded}

\begin{verbatim}
$samsung
    term1    term2 similarity rank
1 samsung     asus  0.7393969    1
2 samsung     acer  0.7289189    2
3 samsung motorola  0.7063525    3
4 samsung    nexus  0.7028567    4
5 samsung    vizio  0.7005945    5
\end{verbatim}

The closest words to ``samsung'' are linked to a series of other brands
that, perhaps, are perceived as being similar or to which it is often compared.
Note that even without data cleaning and pre-processing this is an easy
and straightforward way to detect co-occurrences of competing brands and
brand associations in texts.

We can also visualize how words are placed next to or far from each
other. We can shape the analysis by creating a corpus, pre-processing it,
and creating a document-term matrix. For convenience we have selectred 50 words
(otherwise, if we employ all of them, the image would look too crowded
and it would be impossible to see them). The idea here is to choose just a
few words and then visualize them.

\begin{Shaded}
\begin{Highlighting}[]
\FunctionTok{library}\NormalTok{(tm)}
\FunctionTok{library}\NormalTok{(magrittr)}

\CommentTok{\# Select the Content column}
\NormalTok{reviews }\OtherTok{\textless{}{-}}\NormalTok{ prodreviews}\SpecialCharTok{$}\NormalTok{Content}

\CommentTok{\# Preprocess the corpus}
\NormalTok{corpus }\OtherTok{\textless{}{-}} \FunctionTok{Corpus}\NormalTok{(}\FunctionTok{VectorSource}\NormalTok{(reviews)) }\SpecialCharTok{\%\textgreater{}\%}
  \FunctionTok{tm\_map}\NormalTok{(}\FunctionTok{content\_transformer}\NormalTok{(tolower)) }\SpecialCharTok{\%\textgreater{}\%}
  \FunctionTok{tm\_map}\NormalTok{(removePunctuation) }\SpecialCharTok{\%\textgreater{}\%}
  \FunctionTok{tm\_map}\NormalTok{(removeNumbers) }\SpecialCharTok{\%\textgreater{}\%}
  \FunctionTok{tm\_map}\NormalTok{(removeWords, }\FunctionTok{stopwords}\NormalTok{(}\StringTok{"en"}\NormalTok{)) }\SpecialCharTok{\%\textgreater{}\%}
  \FunctionTok{tm\_map}\NormalTok{(stripWhitespace)}

\CommentTok{\# Extract the words from the document{-}term matrix}
\NormalTok{words }\OtherTok{\textless{}{-}} \FunctionTok{colnames}\NormalTok{(}\FunctionTok{as.matrix}\NormalTok{(}\FunctionTok{DocumentTermMatrix}\NormalTok{(corpus)))}

\CommentTok{\# Remove duplicates and empty strings}
\NormalTok{word\_list }\OtherTok{\textless{}{-}} \FunctionTok{unique}\NormalTok{(words[words }\SpecialCharTok{!=} \StringTok{""}\NormalTok{])}
\CommentTok{\# Select the first 50 words}
\NormalTok{word\_list }\OtherTok{\textless{}{-}} \FunctionTok{head}\NormalTok{(word\_list, }\DecValTok{50}\NormalTok{)}

\CommentTok{\# Check embeddings and remove the empty values, if any}
\NormalTok{embedding }\OtherTok{\textless{}{-}} \FunctionTok{predict}\NormalTok{(model, word\_list, }\AttributeTok{type =} \StringTok{"embedding"}\NormalTok{)}
\NormalTok{embedding }\OtherTok{\textless{}{-}} \FunctionTok{na.omit}\NormalTok{(embedding)}
\end{Highlighting}
\end{Shaded}

We have used the same CBOW model as above, but the result type now is
``embedding'', not ``nearest''. In this case, the 100 words selected
would be represented as series of numerical values. To visualize the
terms, we can use a dimensionality reduction technique which compresses
the high-dimensional spaces into two or three dimensions. In the present
case we can use uniform manifold approximation and projection (UMAP),
available in the \texttt{plotly} library. Several algorithms perform
a similar task, such as t-SNE (Van der Maaten and Hinton 2008).

\begin{Shaded}
\begin{Highlighting}[]
\FunctionTok{library}\NormalTok{(ggrepel)}
\FunctionTok{library}\NormalTok{(plotly)}
\FunctionTok{library}\NormalTok{(umap)}

\CommentTok{\# Visualize embeddings using UMAP}
\NormalTok{vizualization }\OtherTok{\textless{}{-}} \FunctionTok{umap}\NormalTok{(embedding, }\AttributeTok{n\_neighbors =} \DecValTok{15}\NormalTok{, }\AttributeTok{n\_threads =} \DecValTok{2}\NormalTok{)}

\CommentTok{\# Transform into a data frame}
\NormalTok{df }\OtherTok{\textless{}{-}} \FunctionTok{data.frame}\NormalTok{(}\AttributeTok{word =} \FunctionTok{rownames}\NormalTok{(embedding),}
                 \AttributeTok{x =}\NormalTok{ vizualization}\SpecialCharTok{$}\NormalTok{layout[, }\DecValTok{1}\NormalTok{],}
                 \AttributeTok{y =}\NormalTok{ vizualization}\SpecialCharTok{$}\NormalTok{layout[, }\DecValTok{2}\NormalTok{],}
                 \AttributeTok{stringsAsFactors =} \ConstantTok{FALSE}\NormalTok{)}

\CommentTok{\# Plot embeddings using Plotly}
\FunctionTok{plot\_ly}\NormalTok{(df, }\AttributeTok{x =} \SpecialCharTok{\textasciitilde{}}\NormalTok{x, }\AttributeTok{y =} \SpecialCharTok{\textasciitilde{}}\NormalTok{y, }\AttributeTok{type =} \StringTok{"scatter"}\NormalTok{, }\AttributeTok{mode =} \StringTok{\textquotesingle{}text\textquotesingle{}}\NormalTok{, }\AttributeTok{text =} \SpecialCharTok{\textasciitilde{}}\NormalTok{word) }\CommentTok{\#\%\textgreater{}\%}
\end{Highlighting}
\end{Shaded}

\begin{figure}[H]

\centering{

\includegraphics{10.Chapter_files/figure-pdf/fig-w2v2be-1.pdf}

}

\caption{\label{fig-w2v2be}CBOW embeddings visualization}

\end{figure}%

\begin{Shaded}
\begin{Highlighting}[]
  \CommentTok{\#layout(title = "CBOW Embeddings Visualization")}
\end{Highlighting}
\end{Shaded}

The plot displayed in Figure~\ref{fig-w2v2be} shows the distribution of
the words based on their embedding vectors, but transformed into a
two-dimensional space. The closer the words are positioned to each
other, the more related (or similar) they are. If we take a closer look,
we can see that terms like ``hulu'', ``netflix'', as well as
``ebooks'' and ``games'', appear close to each other. This might indicate
that a tablet can serve several purposes, for watching movies, playing games, and
reading ebooks. These words appear as a group, implying that they are
frequently used in similar contexts. Words that are placed far apart,
contextually, have less in common. This word mapping can be used for several purposes in marketing. One is for competitive analysis carried out for two or more brands, another is for recommendation
systems, by discovering which products are frequently mentioned together
in online reviews. Yet, another could be in personalized marketing
campaigns. Last but not least, selecting relevant words for SEO or ads could give a company an advantage over its competition by using the most appropriate
keywords (see Reisenbichler et al. 2022).

 \texttt{proxy} library.

\begin{Shaded}
\begin{Highlighting}[]
\FunctionTok{library}\NormalTok{(proxy)}

\CommentTok{\# Use the proxy package to compute the cosine distance}
\NormalTok{embedding\_cos }\OtherTok{\textless{}{-}} \FunctionTok{dist}\NormalTok{(embedding, }\AttributeTok{method =} \StringTok{"cosine"}\NormalTok{)}

\CommentTok{\# Perform hierarchical clustering, plot the dendrogram}
\FunctionTok{plot}\NormalTok{(}\FunctionTok{as.dendrogram}\NormalTok{(}\FunctionTok{hclust}\NormalTok{(embedding\_cos)), }\AttributeTok{horiz =} \ConstantTok{FALSE}\NormalTok{, }\AttributeTok{cex =} \FloatTok{0.2}\NormalTok{)}
\end{Highlighting}
\end{Shaded}

\begin{figure}[H]

\centering{

\includegraphics{10.Chapter_files/figure-pdf/fig-cosine-1.pdf}

}

\caption{\label{fig-cosine}Cluster dendrogram of embedded words}

\end{figure}%

As in Figure~\ref{fig-cosine} above, we can see from looking at the
dendrogram that words like ``hulu'' and ``netflix'' appear to be closely
related, as is ``games''. Obviously, the same device
can be used to watch the two streaming services as well as to play
games.

One of the most common things we can do with vectors is to perform
arithmetical operations. What vectors might we create if we take the
embedding for ``kindle'' and move away from it by subtracting ``good'', and
thus reducing its positive association? The result given by
adding ``bad'', would lean towards terms that are not favorable. We
could think of it in this way: ``If (the word) kindle is associated with good, what
words associate kindle with bad?'' This example demonstrates how we can
 process word vectors arithmetically, and the outcome shows how useful this
operation can be.

\begin{Shaded}
\begin{Highlighting}[]
\CommentTok{\# Get the vectors for specific words: \textquotesingle{}amazon\textquotesingle{}, \textquotesingle{}apple\textquotesingle{}, \textquotesingle{}good\textquotesingle{}, \textquotesingle{}bad\textquotesingle{}, \textquotesingle{}screen\textquotesingle{}}
\NormalTok{kindle }\OtherTok{\textless{}{-}} \FunctionTok{predict}\NormalTok{(model, }\AttributeTok{newdata =} \StringTok{"kindle"}\NormalTok{, }\AttributeTok{type =} \StringTok{"embedding"}\NormalTok{)}
\NormalTok{good }\OtherTok{\textless{}{-}} \FunctionTok{predict}\NormalTok{(model, }\AttributeTok{newdata =} \StringTok{"good"}\NormalTok{, }\AttributeTok{type =} \StringTok{"embedding"}\NormalTok{)}
\NormalTok{bad }\OtherTok{\textless{}{-}} \FunctionTok{predict}\NormalTok{(model, }\AttributeTok{newdata =} \StringTok{"bad"}\NormalTok{, }\AttributeTok{type =} \StringTok{"embedding"}\NormalTok{)}

\CommentTok{\# Perform vector arithmetic for the word kindle}
\NormalTok{result\_kindle }\OtherTok{\textless{}{-}}\NormalTok{ kindle }\SpecialCharTok{{-}}\NormalTok{ good }\SpecialCharTok{+}\NormalTok{ bad}

\CommentTok{\# Find the top 10 words closest to the result vector for kindle}
\NormalTok{similar\_kindle }\OtherTok{\textless{}{-}} \FunctionTok{predict}\NormalTok{(model, }\AttributeTok{newdata =}\NormalTok{ result\_kindle, }\AttributeTok{type =} \StringTok{"nearest"}\NormalTok{, }\AttributeTok{top\_n =} \DecValTok{10}\NormalTok{)}
\FunctionTok{print}\NormalTok{(similar\_kindle)}
\end{Highlighting}
\end{Shaded}

\begin{verbatim}
$kindle
      term similarity rank
1   kindle  0.9984323    1
2   kindel  0.7714513    2
3  kindles  0.7541807    3
4    kidle  0.7478809    4
5   kinkle  0.7471732    5
6    kndle  0.7398569    6
7   kindal  0.7393759    7
8   kendle  0.7232242    8
9    kinle  0.7011244    9
10       k  0.6855328   10
\end{verbatim}

This unexpected result comes from the series of words which are mistypes
of the word ``kindle'' and shows that careful pre-processing is essential
for meaningful text analysis.

Let us perform further operations: We can add positive characteristics,
such as ''good'', ''fast'', and ''cheap'', and remove words like ''bad'', ''slow'', and
''expensive''. In doing so, we can try to shift towards more positive qualities (the one we have added) and move away from negative ones (the ones we have subtracted). After
this, we can check which words are embedded close to both ''kindle'' and ''ipad''.

\begin{Shaded}
\begin{Highlighting}[]
\CommentTok{\# Get the vectors for the kindle and ipad}
\NormalTok{kindle }\OtherTok{\textless{}{-}} \FunctionTok{predict}\NormalTok{(model, }\AttributeTok{newdata =} \StringTok{"kindle"}\NormalTok{, }\AttributeTok{type =} \StringTok{"embedding"}\NormalTok{)}
\NormalTok{ipad }\OtherTok{\textless{}{-}} \FunctionTok{predict}\NormalTok{(model, }\AttributeTok{newdata =} \StringTok{"ipad"}\NormalTok{, }\AttributeTok{type =} \StringTok{"embedding"}\NormalTok{)}

\CommentTok{\# Get the vectors for good, bad, fast, slow, cheap and expensive}
\NormalTok{good }\OtherTok{\textless{}{-}} \FunctionTok{predict}\NormalTok{(model, }\AttributeTok{newdata =} \StringTok{"good"}\NormalTok{, }\AttributeTok{type =} \StringTok{"embedding"}\NormalTok{)}
\NormalTok{bad }\OtherTok{\textless{}{-}} \FunctionTok{predict}\NormalTok{(model, }\AttributeTok{newdata =} \StringTok{"bad"}\NormalTok{, }\AttributeTok{type =} \StringTok{"embedding"}\NormalTok{)}
\NormalTok{fast }\OtherTok{\textless{}{-}} \FunctionTok{predict}\NormalTok{(model, }\AttributeTok{newdata =} \StringTok{"fast"}\NormalTok{, }\AttributeTok{type =} \StringTok{"embedding"}\NormalTok{)}
\NormalTok{slow }\OtherTok{\textless{}{-}} \FunctionTok{predict}\NormalTok{(model, }\AttributeTok{newdata =} \StringTok{"slow"}\NormalTok{, }\AttributeTok{type =} \StringTok{"embedding"}\NormalTok{)}
\NormalTok{cheap }\OtherTok{\textless{}{-}} \FunctionTok{predict}\NormalTok{(model, }\AttributeTok{newdata =} \StringTok{"cheap"}\NormalTok{, }\AttributeTok{type =} \StringTok{"embedding"}\NormalTok{)}
\NormalTok{expensive }\OtherTok{\textless{}{-}} \FunctionTok{predict}\NormalTok{(model, }\AttributeTok{newdata =} \StringTok{"expensive"}\NormalTok{, }\AttributeTok{type =} \StringTok{"embedding"}\NormalTok{)}

\CommentTok{\# Perform vector arithmetic for the kindle word}
\NormalTok{result\_kindle }\OtherTok{\textless{}{-}}\NormalTok{ kindle }\SpecialCharTok{+}\NormalTok{ good }\SpecialCharTok{+}\NormalTok{ fast }\SpecialCharTok{+}\NormalTok{ cheap }\SpecialCharTok{{-}}\NormalTok{ bad }\SpecialCharTok{{-}}\NormalTok{ slow }\SpecialCharTok{{-}}\NormalTok{ expensive}

\CommentTok{\# Perform vector arithmetic for the ipad word}
\NormalTok{result\_ipad }\OtherTok{\textless{}{-}}\NormalTok{ ipad }\SpecialCharTok{+}\NormalTok{ good }\SpecialCharTok{+}\NormalTok{ fast }\SpecialCharTok{+}\NormalTok{ expensive }\SpecialCharTok{{-}}\NormalTok{ bad }\SpecialCharTok{{-}}\NormalTok{ slow }\SpecialCharTok{{-}}\NormalTok{ cheap}

\CommentTok{\# Find the top 10 words closest to each result vector}
\NormalTok{similar\_kindle }\OtherTok{\textless{}{-}} \FunctionTok{predict}\NormalTok{(model, }\AttributeTok{newdata =}\NormalTok{ result\_kindle, }\AttributeTok{type =} \StringTok{"nearest"}\NormalTok{, }\AttributeTok{top\_n =} \DecValTok{10}\NormalTok{)}
\NormalTok{similar\_ipad }\OtherTok{\textless{}{-}} \FunctionTok{predict}\NormalTok{(model, }\AttributeTok{newdata =}\NormalTok{ result\_ipad, }\AttributeTok{type =} \StringTok{"nearest"}\NormalTok{, }\AttributeTok{top\_n =} \DecValTok{10}\NormalTok{)}

\CommentTok{\# Print the results for kindle}
\FunctionTok{print}\NormalTok{(similar\_kindle)}
\end{Highlighting}
\end{Shaded}

\begin{verbatim}
$kindle
          term similarity rank
1       kindle  0.9957064    1
2         fast  0.8844871    2
3        quick  0.8438365    3
4        sleek  0.8350880    4
5  wonderfully  0.8345558    5
6         good  0.8332866    6
7       kindel  0.8235799    7
8       simple  0.8204671    8
9      solidly  0.8192567    9
10        well  0.8092886   10
\end{verbatim}

The words which show the positive qualities of Amazon's Kindle are related to the response time of the machine (fast, quick),
that it is sleek, perhaps that it does some task extremely well, that its use is
straightforward, and the tablet looks solid, among others.

\begin{Shaded}
\begin{Highlighting}[]
\CommentTok{\# Print the results for ipad}
\FunctionTok{print}\NormalTok{(similar\_ipad)}
\end{Highlighting}
\end{Shaded}

\begin{verbatim}
$ipad
         term similarity rank
1  economical  0.9407232    1
2     ereader  0.9154697    2
3         pad  0.9062951    3
4   expensive  0.8988186    4
5        imac  0.8940031    5
6      iphone  0.8841532    6
7      itouch  0.8775398    7
8     elegant  0.8763888    8
9       fluid  0.8541148    9
10  efficient  0.8540065   10
\end{verbatim}

The positive words related to Apple's iPad show a degree of ambivalence,
since it seems to be perceived as economical by some, but expensive
by others. Also, the iPad is associated with an e-reader and that it is
integrated into Apple's ecosystem (e.g., ''imac'', ''iphone'', ''itouch''). Other
positive associations include ''efficient'' and ''elegant'', and perhaps that images
are ''fluid'' on the screen.

Marketers can use such information in several ways. For example,
for Kindle, Amazon could emphasize the simplicity, yet solidity of the
tablet, as well as its quick response to commands. Apple's marketing
could focus on its design and its integration with other Apple devices,
emphasizing that the iPad is not just a tablet but an efficient tool.

\section{Transformers in action}\label{transformers-in-action}

In this section, we will illustrate the use of transformers through the R
package text (Kjell, Giorgi, and Schwartz 2023). This package
provides an interface with trained deep-learning models which allow
several NLP tasks to use the power of transformers and Python
via R. Several tasks can be carried out using transformers, for
example topic modeling, named-entity recognition, sentiment analysis,
and many more. These have already been explored in previous
chapters using more classical approaches. In this section we will
focus on two tasks, one is text summarizing, developing
what we discussed in
Chapter~\ref{sec-named-entity-recognition-and-summarization} further, and the
second is related to searching for information  (i.e., questions and
answers).

\subsection{Text Summarizing}\label{text-summarization}

Summarizing is important in marketing (Hartmann and Netzer 2023)
because it helps marketing analysts condense large volumes of
customer feedback into shorter, concise insights which can be used to inform further action. It is useful for understanding customer sentiment, or identifying
trends at scale. In the next example we summarize a series of
randomly selected customer reviews. We have set a seed for reproducibility, filtered the dataset for the Amazon brand, and randomly selected 100
reviews to speed up the computational process. These reviews
are then linked into a single string to produce a single
document. Then, we applied the \texttt{textSum()} function and by using a
trained deep-learning model (t5-small) we have summarized the text. The
summary is generated as a string that is between 5 and
100 tokens long. The result is stored in the \texttt{ts}
variable which can be printed using the \texttt{cat()} function.

\begin{Shaded}
\begin{Highlighting}[]
\FunctionTok{library}\NormalTok{(text)}

\CommentTok{\# Install text required python packages in a conda environment (with defaults).}
\CommentTok{\# textrpp\_initialize(save\_profile = TRUE) saves the profile}
\FunctionTok{textrpp\_install}\NormalTok{()}
\FunctionTok{textrpp\_initialize}\NormalTok{()}

\CommentTok{\# Select the Amazon brand}
\NormalTok{amazon\_reviews }\OtherTok{\textless{}{-}}\NormalTok{ prodreviews[prodreviews}\SpecialCharTok{$}\NormalTok{Brand }\SpecialCharTok{==} \StringTok{"Amazon"}\NormalTok{, ]}


\CommentTok{\# Set seed for reproducibility}
\FunctionTok{set.seed}\NormalTok{(}\DecValTok{123}\NormalTok{) }

\CommentTok{\# Sample 1000 indices from the available reviews}
\NormalTok{sampled\_ind }\OtherTok{\textless{}{-}} \FunctionTok{sample}\NormalTok{(}\DecValTok{1}\SpecialCharTok{:}\FunctionTok{nrow}\NormalTok{(amazon\_reviews), }\AttributeTok{size =} \DecValTok{100}\NormalTok{, }\AttributeTok{replace =} \ConstantTok{FALSE}\NormalTok{)}

\CommentTok{\# Extract the content of the sampled reviews}
\NormalTok{sampled\_reviews }\OtherTok{\textless{}{-}}\NormalTok{ amazon\_reviews}\SpecialCharTok{$}\NormalTok{Content[sampled\_ind]}

\CommentTok{\# Concatenate the text into one big string}
\NormalTok{big\_review }\OtherTok{\textless{}{-}} \FunctionTok{paste}\NormalTok{(sampled\_reviews, }\AttributeTok{collapse =} \StringTok{" "}\NormalTok{)}

\CommentTok{\# Perform the summarization operation, use the t5{-}small model}
\NormalTok{ts }\OtherTok{\textless{}{-}} \FunctionTok{textSum}\NormalTok{(big\_review, }\AttributeTok{min\_length =} \DecValTok{5}\NormalTok{L, }\AttributeTok{max\_length =} \DecValTok{100}\NormalTok{L, }\AttributeTok{model =} \StringTok{"t5{-}small"}\NormalTok{, }\AttributeTok{device =} \StringTok{"gpu"}\NormalTok{, }\AttributeTok{tokenizer\_parallelism =} \ConstantTok{TRUE}\NormalTok{, }\AttributeTok{logging\_level =} \StringTok{"warning"}\NormalTok{, }\AttributeTok{return\_text =} \ConstantTok{TRUE}\NormalTok{, }\AttributeTok{return\_tensors =} \ConstantTok{FALSE}\NormalTok{, }\AttributeTok{clean\_up\_tokenization\_spaces =} \ConstantTok{FALSE}\NormalTok{, }\AttributeTok{set\_seed =} \DecValTok{42}\NormalTok{L )}
\end{Highlighting}
\end{Shaded}

\begin{Shaded}
\begin{Highlighting}[]
\CommentTok{\# View the results}
\FunctionTok{cat}\NormalTok{(ts[[}\DecValTok{1}\NormalTok{,}\DecValTok{1}\NormalTok{]])}
\end{Highlighting}
\end{Shaded}

\begin{verbatim}
the Kindle Fire is a great device for reading and watching movies . 
it's a good way to get a Kindle, but it isn't as good as the iPad. 
if you're a fan of the e-book, you'll be able to use it for a lot of time.
\end{verbatim}

You can see the good and bad points of the device immediately. It is also possible to select different data, for example, by choosing another brand or making a
comparison between them. The more data you have, the more precise the
summaries will be.

\subsection{Information searches, questions and answering using
transformers
(Q\&A)}\label{information-search-question-and-answering-using-transformers-qa}

Searching for information and customer feedback plays a central role in modern
marketing. Transformers have revolutionized this by providing
accurate, context-aware retrieval systems, allowing brands to analyze vast
amounts of consumer queries and intent with high precision. Also
transformer-powered chatbots have had an impact on customer interactions by
automating the question-and-answer process
(Smith and Albaum 2010). Q\&As give direct insights into customer
experiences and pain points. In contrast to surveys, Q\&As
provide real-time feedback, providing unexpected insights, a
more conversational approach to dynamic exchanges, and
scalability. Q\&As will not replace surveys, especially when there is
the need for structured information on aspects of a product with a
corresponding statistical analysis. Nevertheless, there is a growing
body of literature which is exploring the ability of transformer-based
language models to mimic consumer research (for an overview, see N.
Arora, Chakraborty, and Nishimura 2025; Sarstedt et al. 2024).
Despite the enthusiasm for ``silicon sampling'' there is recent research
which has also identified its limitations (e.g., Goli and Singh 2024;
Stromberg et al. 2025).

To illustrate this, here are a list of questions that we would like to see
answered.

\begin{Shaded}
\begin{Highlighting}[]
\CommentTok{\# List of questions to be answered}
\NormalTok{questions }\OtherTok{\textless{}{-}} \FunctionTok{c}\NormalTok{(}
\StringTok{"How is the Kindle screen?"}\NormalTok{,}
\StringTok{"How is the Amazon customer service?"}\NormalTok{,}
\StringTok{"How much memory does the Kindle Fire have?"}\NormalTok{,}
\StringTok{"Where is the Kindle tablet camera?"}\NormalTok{,}
\StringTok{"Does the brand Amazon offer good value for the money?"}
\NormalTok{)}
\end{Highlighting}
\end{Shaded}

We have stored the questions in a vector, then iterated each question by
applying the \texttt{textQA()} function, dedicated to asking questions. There are several models that can be
used, you can try substituting the one proposed here with different
ones, check the \emph{text} package specifications for this. Also, the answers can vary depending on the exact prompt used to question the
language model (see P. Liu et al., 2023, on prompting and the prompt
engineering of language models).

\begin{Shaded}
\begin{Highlighting}[]
\CommentTok{\# Initialize a list to store the answers}
\NormalTok{answers }\OtherTok{\textless{}{-}} \FunctionTok{vector}\NormalTok{(}\StringTok{"list"}\NormalTok{, }\FunctionTok{length}\NormalTok{(questions))}

\CommentTok{\# Iterate over each question and apply the textQA function}
\ControlFlowTok{for}\NormalTok{ (i }\ControlFlowTok{in} \FunctionTok{seq\_along}\NormalTok{(questions)) \{}
\NormalTok{  answers[[i]] }\OtherTok{\textless{}{-}} \FunctionTok{textQA}\NormalTok{(}
    \AttributeTok{question =}\NormalTok{ questions[i],}
    \AttributeTok{context =}\NormalTok{ big\_review,}
    \AttributeTok{model =} \StringTok{"deepset/roberta{-}base{-}squad2"}\NormalTok{,}
    \AttributeTok{device =} \StringTok{"cpu"}\NormalTok{,}
    \AttributeTok{tokenizer\_parallelism =} \ConstantTok{FALSE}\NormalTok{,}
    \AttributeTok{logging\_level =} \StringTok{"warning"}\NormalTok{,}
    \AttributeTok{top\_k =} \DecValTok{10}\NormalTok{L,}
    \AttributeTok{doc\_stride =} \DecValTok{128}\NormalTok{L,}
    \AttributeTok{max\_answer\_len =} \DecValTok{20}\NormalTok{L,}
    \AttributeTok{max\_seq\_len =} \DecValTok{384}\NormalTok{L,}
    \AttributeTok{max\_question\_len =} \DecValTok{64}\NormalTok{L,}
    \AttributeTok{handle\_impossible\_answer =} \ConstantTok{FALSE}\NormalTok{,}
    \AttributeTok{set\_seed =} \DecValTok{202208}\NormalTok{L}
\NormalTok{  )}
\NormalTok{\}}
\end{Highlighting}
\end{Shaded}

\begin{Shaded}
\begin{Highlighting}[]
\CommentTok{\# Print the answers}
\ControlFlowTok{for}\NormalTok{ (i }\ControlFlowTok{in} \FunctionTok{seq\_along}\NormalTok{(questions)) \{}
  \FunctionTok{cat}\NormalTok{(}\StringTok{"Question:"}\NormalTok{, questions[i], }\StringTok{"}\SpecialCharTok{\textbackslash{}n}\StringTok{Answer:"}\NormalTok{, }\FunctionTok{paste}\NormalTok{(}\FunctionTok{unlist}\NormalTok{(answers[[i]][}\DecValTok{4}\NormalTok{]), }\AttributeTok{collapse =} \StringTok{", "}\NormalTok{), }\StringTok{"}\SpecialCharTok{\textbackslash{}n\textbackslash{}n}\StringTok{"}\NormalTok{)}
\NormalTok{\}}
\end{Highlighting}
\end{Shaded}

\begin{verbatim}
Question: How is the Kindle screen? 
Answer: back-lit, great, great, back-lit, perfectly  lit and easy to 
navigate or read, 10.1, glare, great, great screen display, sharp
and videos are very crisp 

Question: How is the Amazon customer service? 
Answer: reat, great, you honestly have a customer for life, great as
always, great as always, loyal, expedited shipping, honestly have a 
customer for life, monopolizing these, expedited shipping 

Question: How much memory does the Kindle Fire have? 
Answer: 8 GB, 8gig, 8gig, 8 GB memory, 7 inch, ten, No extra memory
slot, 7-inch, 7"", No extra memory slot? what for? use the cloud 

Question: Where is the Kindle tablet camera? 
Answer: front & back, GPS, front & back camera, there isn't a camera, 
kindle fire, sd slot, smartphone, It's almost an IPad, I-pad, laptop 

Question: Does the brand Amazon offer good value for the money? 
Answer: this technology could be so affordable, its well worth the price,
you can't go wrong with this at all, Amazon has a winner, the Kindle Fire
is perfectfor functionality and price, it is so worth the money, well worth
the money!!!, Kindle is a bargain, Amazon has a winner!, it's much cheaper 
\end{verbatim}

The answers seem to provide useful information about the
characteristics of the Kindle Fire, or about the Amazon brand in general. The precision of the
answer, allows for further questions or further analysis. You could also select
a different product and ask the same questions for a comparative
analysis.

Throughout this book, we have come across the many standard and modern
approaches for tackling text analysis in marketing. The ever-changing
technological support in terms of packages equips the marketing analyst
with powerful tools for analyzing and discovering latent messages in user-generated content by transforming huge volumes of text into manageable parcels of information.

As an end note, the models presented in this chapter find text, or
generate text, by repeating (like a parrot) patterns found in the
training data, but truly understanding or reasoning about language is
still a human quality. Machines can produce output based on statistical
probabilities (the word stochastic is used in these cases), but do
not (yet) have a comprehension of its meaning (Bender et
al. 2021).

\emph{Exercise starts}

\begin{enumerate}
\def\labelenumi{\arabic{enumi}.}
\tightlist
\item
Practise using other vectors, such as ''kindle'' - ''good'' + ''ipad'', and find out
  things about ipad.
\item
  Instead of CBOW, use the skipgram method to train a word-embedding
  data set.
\item
  Use the function \texttt{textNER()} from the \texttt{text} package for
  named-entity recognition tasks.
\item
  Use the function \texttt{textTopics()} from the \texttt{text} package
  to extract topics.
\item
  Select a positive and negative subset from the main data set, summarize
  each of them and compare the results.
\end{enumerate}

\emph{Exercise ends}

\emph{Conclusion starts}

The chapter began with an introduction to word embedding and the
Word2Vec algorithm. It presented a walk-through the
transformation of a corpus into a vector representation of the words.
The first example started with a two dimensional array which showed the
way to construct to a manageable structure. From this point on, multidimensional
vectors were easier to be understand. Word2Vec was analyzed in detail, and
both continuous bag-of-words and skip-gram were mentioned as techniques, as well as
negative sampling. A simple practical example was offered for the better
understanding of the training mechanism for neural networks.
An example of backpropagation showed how training
reaches the optimal values of embedding. The next step was the explanation
of transformers, an AI algorithm which has received a lot of attention. There was also a succinct explanation of the use of encoders,
encoders-decoders, and decoders. Examples of word embedding
were given to show their applications in marketing, tackling
analogies, clustering, and also for vector operations. The chapter ended with
the presentation of two operations using transformers via the
\texttt{text}package. The first dealt with summarizing, the second with the implementation of a Q\&A mechanism.

\emph{Conclusion ends}